\documentclass[10pt,a4paper]{amsart}
\usepackage[margin=19mm]{geometry}
\usepackage{amsmath,amssymb,mathtools}
\usepackage[scr=rsfs,cal=euler]{mathalfa}
\usepackage{xfrac}
\usepackage[unicode,hidelinks,bookmarks=false]{hyperref}
\newcommand{\ie}{i.e.\ }
\newcommand{\cf}{cf.\ }
\newcommand{\resp}{respectively\ }
\newcommand{\Subsection}{\subsection}
\providecommand{\nobreakdash}{\nobreak\hbox{-}}
\setlength{\emergencystretch}{2em}
\multlinegap0pt
\usepackage{bm}
\usepackage{bbm}
\usepackage{dsfont}
\usepackage{xspace}
\usepackage{cancel}
\usepackage{mathtools}
\usepackage{amsthm}
\makeatletter
\@ifundefined{theorem}{\newtheorem{theorem}{Theorem}}{}
\@ifundefined{lemma}{\newtheorem{lemma}[theorem]{Lemma}}{}
\makeatother
\newcommand{\N}{\mathbb{N}}
\newcommand{\R}{\mathbb{R}}
\newcommand{\dx}[1]{\,\mathrm{d}#1}  % Proper differential operator
\newcommand{\eps}{\varepsilon}
\title[Sharp mean-field estimates for the repulsive log gas in any ]{Sharp mean-field estimates for the repulsive log gas in any dimension}
\author{Matias G. Delgadino, Rishabh S. Gvalani}
\date{Source: arXiv:2506.22083v1 (CC BY 4.0)}
\begin{document}
\maketitle

\noindent\emph{Source and licence.} Sharp mean-field estimates for the repulsive log gas in any dimension. Version arXiv:2506.22083v1 (CC BY 4.0), distributed under the Creative Commons Attribution 4.0 International licence, as recorded in the arXiv licence metadata for this exact version and shown on the abstract page https://arxiv.org/abs/2506.22083v1. Full rights notice: licenses/arxiv-2506.22083-CC-BY-4.0.txt.\par\smallskip

\noindent\textit{Editorial scope.} Counted unit: the Lemma whose source label is \texttt{lem:corineq} in the article (source0.tex lines 432--444, source label \texttt{lem:corineq}), delivered together with the article's own complete original proof of it (source0.tex lines 445--571, one \texttt{proof} environment of 127 lines). No mathematics is written, repaired or re-derived: statement and proof are the article's own text line for line. Required context retained uncounted: none. Every definition, display and label of those lines is retained unchanged. Omitted from this package and not redistributed: 1--431, 572--820. Nothing inside a retained excerpt is omitted or altered: every retained body is delivered exactly as the article's lines are written, so no omission receipt of any kind accompanies this package. Citations: the retained excerpts carry no citation command at all, so no bibliography entry of the article is transcribed and no reference is redistributed. Reference adaptation: none was needed and none was made; every reference inside the delivered unit resolves inside it, so no unresolved reference or citation is produced. Printed slips kept literal: no printed word, symbol, hypothesis or number of the counted statement or of its proof was altered. Layout and class: the article's own class is \texttt{amsart} with packages including amsmath, amssymb, amsthm, hyperref, mathtools, esint, enumitem, accents, xcolor, geometry, tikz-cd, cases, cleveref, autonum; the wrapper is the lane's pinned \texttt{amsart} editorial wrapper at 10pt on A4, which restates the theorem-like environments the retained text needs and adds the article's own macro definitions that the retained text needs (\texttt{\textbackslash N}, \texttt{\textbackslash R}, \texttt{\textbackslash dx}, \texttt{\textbackslash eps}), with extra wrapper packages (bm, bbm, dsfont, xspace, cancel, mathtools). The delivered numbering is the unit's own. Assets: no retained span contains a figure environment, a graphics-inclusion command, a PDF-page inclusion, a table or a TikZ picture; no image file is copied into this package, the package declares no asset dependency and no image file is redistributed with it. Licence: version arXiv:2506.22083v1 of this article is distributed under the Creative Commons Attribution 4.0 International licence, as recorded in the arXiv licence metadata for this exact version and shown on the abstract page https://arxiv.org/abs/2506.22083v1. A full rights notice is delivered at licenses/arxiv-2506.22083-CC-BY-4.0.txt. Characters: every retained excerpt is pure ASCII, so the delivered engine, XeLaTeX with its default Latin Modern text font, reports no missing character.
\begin{lemma}[Correlation inequality]\label{lem:corineq}
Fix $p \in \N$ and $\gamma \in [1/2,1)$, such that $1\le p\le N$. Let $G:\Omega^2 \to \R$ be such that $G(x,y)=G(y,x)$ for all $x,y \in \Omega$,
\begin{equation}
  \sup_{x\in \Omega}\lVert G(x,\cdot) \rVert_{L^p(\Omega;\bar\rho)} <\infty\,, \quad \textrm{ and }\quad \int_{\Omega}G(x,y) \, \dx{\bar\rho}(y)=0\,,\, \forall x \in \Omega\, .
\end{equation}

Then, there exists a constant $C_p<\infty$, such that
\begin{align}
&\mathbb{{E}}\left[\left|\frac{1}{N}\sum_{i\neq j}^{N}G(X_{i},X_{j})\right|^{p}\right]\\&\qquad \leq C_p\left(\frac{1}{N^{p-1-\lfloor \gamma p \rfloor}} \sup_{x \in \Omega} \lVert  G(x,\cdot)\rVert_{L^p(\Omega;\bar\rho)}^{p}\, +  \sup_{x \in \Omega} \lVert  G(x,\cdot)\rVert_{L^{2(p-\lceil \gamma p \rceil)}(\Omega;\bar\rho)}^{p} \right)\, ,
\label{eq:corineq}
\end{align}
where $(X_i)_{i=1,\dots,N}$ are independent $\bar\rho$-distributed random variables.
\end{lemma}

\begin{proof}
  Before we can proceed to the main body of the proof, we need to introduce some notation to keep track of the various terms that show up when we expand the term on the left hand side of~\eqref{eq:corineq}. Consider the set  
\begin{equation}
\label{eq:defI}
A:= \{(i,j)\in \{1,\cdots,N\}^2:i\neq j\}\, .
\end{equation}
We call a map $I_p:A \to \N\cup\{0\}$ a $p$-multiindex if 
\begin{align}
\sum_{(i,j)\in A}I_p((i,j))=p\, .
\end{align}
Given a $p$-multiindex $I_p$, we associate to it the following multiplicity vector $M_{I_p}= (m_{1,I_p},\dots,m_{N,I_p})\in (\N\cup\{0\})^N$, where
\begin{equation}
m_{i,I_p}:= \sum_{j =1}^N \left(I_p((i,j)) + I_p((j,i))\right)\, .
\end{equation}
We will also need to use the following restricted set of $p$-multiindices
\begin{equation}
\label{eq:restrictedset}
\mathcal{E}_{p}:= \{I_p \text{ is a }p \text{-multiindex}:m_{i,I_p}\neq 1,\, \forall i \in \{1,\dots,N\}\}\,.
\end{equation}
Given a $p$-multiindex $I_p$, we count how many active variables the multiindex has
\begin{equation}
\mathrm{act}(I_p):=\left|B_{I_p}\right|, \qquad B_{I_p}:=\{i \in {1, \dots,N}:m_{i,I_p}\neq 0\}\,.
\end{equation}
We further refine $\mathcal{E}_p$ by considering the amount of active variables
\begin{equation}
\mathcal{E}_{p,\ell}:= \{I_p \in \mathcal{E}_{p}: \mathrm{act}(I_p)=\ell\}\,,
\end{equation}
and rewrite it into pairwise disjoint union
\begin{equation}
\mathcal{E}_p= \bigcup_{\ell=1}^p \mathcal{E}_{p,\ell}\, ,
\end{equation}
where we have used that if $I_p \in \mathcal{E}_p$ then we can must have  at most $p$ active variables that $\mathrm{act}(I_p) \leq p$. For $I_p\in\mathcal{E}_{p,\ell}$, for any $i_0\in \{1,\dots,N\}$ we can obtain the bound
\begin{equation}\label{boundonmi}
    m_{i_0,I_{p}}\le 2p-2(\ell-1),
\end{equation}
this follows from the bound
$$
m_{i_0,I_{p}}+2(\ell-1) \le \sum_{i=1}^Nm_{i,I_{p}}=\sum_{(i,j)\in A}I_p((i,j))+I_p((j,i))=2p,
$$
where we use that for any active particle $m_{i,I_{p}}\ge 2$ by definition of $I_p\in\mathcal{E}_p$. Finally, we state the following estimate on the cardinality of each set
\begin{equation}
|\mathcal{E}_{p,\ell}|\le  {N\choose\ell}(\ell^2 -\ell)^p\le C_p N^{\ell}
\label{eq:sizeE}\, .
\end{equation}


With these concepts in mind, we proceed to estimate the quantity in the statement of the lemma. We have
\begin{align}
 \, \mathbb{{E}}\left[\left|\frac{1}{N}\sum_{i\neq j}^{N}G(X_{i},X_{j})\right|^{p}\right]
=&\, \frac{{1}}{N^{p}}\int_{\Omega^N}\sum_{I_{p}}\frac{p!}{\prod_{(i,j)\in A}I_p(i,j)! }\prod_{(i,j)\in A}G(x_{i,},x_{j})^{I_p((i,j))}\,\mathrm{{d}}\bar\rho^{\otimes N} \\
=&\, \frac{{1}}{N^{p}}\int_{\Omega^N}\sum_{I_{p}\in\mathcal{{E}}_{p}}\frac{p!}{\prod_{(i,j)\in A}I_p(i,j)! } \prod_{(i,j)\in A}G(x_{i,},x_{j})^{I_p((i,j))}\,\mathrm{{d}}\bar\rho^{\otimes N} \\
\leq & \, \frac{C_p}{N^{p}}\int_{\Omega^N}\sum_{\ell=2}^p\sum_{I_{p}\in\mathcal{{E}}_{p,\ell}} \prod_{(i,j)\in A}|G(x_{i,},x_{j})|^{I_p((i,j))}\,\mathrm{{d}}\bar\rho^{\otimes N} \, .\label{eq:termtosplit}
\end{align}
Note that the restriction $I_p\in\mathcal{E}_p$ in the second equality stems from the fact that 
\begin{equation}
\int_{\Omega}G(x,y)\:\mathrm{{d}\bar\rho}(x)=\int_{\Omega}G(x,y)\:\mathrm{{d}\bar\rho}(y)=0 \, ,
\end{equation}
from which it follows that the corresponding terms in the sum are zero. Given $I_{p} \in \mathcal{E}_{p,\ell}$, enumerating the elements of $B_{I_p}$ as $\{i_1,\dots,i_\ell\}$, we define the set of active indices that contain $i_k$ by
\begin{equation}
A_{k,I_p}:= \left\{ (i,j) \in A: (i =i_k \;\mbox{or}\; j=i_k)\;\mbox{and}\; I_p(i,j)>0 \right\} \, \qquad\mbox{for $k=1,\dots,\ell$.}
\end{equation} 
For any $I_p \in \mathcal{E}_{p,\ell}$ we can integrate out the $N-\ell$ variables that do not lie in $B_{I_p}$. We then decompose the remaining variables into a disjoint union of the sets $C_{k,I_p}\coloneqq A_{k,I_p}\setminus \bigcup_{m=1}^{k-1}A_{m,I_p}$ for $k=1,\dots,\ell-1$ to derive the identity
\begin{align}
&\, \int_{\Omega^N}\prod_{(i,j)\in A}|G(x_{i,},x_{j})|^{I_p((i,j))}\,\mathrm{{d}}\bar\rho^{\otimes N}\\
= & \, \int_{\Omega^{\ell}}\prod_{k=1 }^{\ell-1}\left(\prod_{(i,j)\in C_{k,I_p}}|G(x_{i,},x_{j})|^{I_p((i,j))}\right)\mathrm{{d}}\bar\rho^{\otimes \ell}(x_{i_1},...,x_{i_\ell})\label{eq:breakingupA}
\end{align}

To estimate \eqref{eq:breakingupA}, we proceed inductively for each element in $k=1,...,\ell-1$. We first introduce some notation and bounds.
\begin{align}
\gamma_k:=& \,\sum_{(i,j)\in C_{k,I_p}}I_p((i,j))\\\leq&\,\sum_{(i,j)\in A_{k,I_p}}I_p((i,j))\\
\eqqcolon&\, m_{k,p}\le \begin{cases}
    p&\ell<\frac{p}{2}+1\\
    2p-2(\ell-1)& \ell\ge \frac{p}{2}+1
\end{cases}
\, ,
\end{align}
where for the last inequality we have applied the bounds from \eqref{boundonmi}.
We start by bounding the integrand that only depends on $x_{i_1}$, and apply H\"older's inequality with exponents $p_j=\gamma_1/I_p(i_1,j)$ to obtain
\begin{align}
&\int_{\Omega}\prod_{(i,j)\in C_{1,I_p}}|G(x_{i,},x_{j})|^{I_p((i,j))}\mathrm{{d}}\bar\rho(x_{i_1})\\
\le & \, \prod_{(i,j)\in C_{1,I_p}} \left(\int_\Omega |G(x_{i_1},x_{j})|^{\gamma_1}\;\dx \bar\rho(x_{i_1})\right)^{I_p(i_1,j)/\alpha_1}\\
\le & \, \prod_{(i,j)\in C_{1,I_p}}\sup_x \left(\int_\Omega |G(x_{i_1},x)|^{\gamma_1}\;\dx \bar\rho(x_{i_1})\right)^{I_p(i_1,j)/\alpha_1}\\
= & \,\sup_x\int_\Omega |G(x_{i_1},x)|^{\gamma_1}\;\dx \bar\rho(x_{i_1})\\
\leq &   \,
 \begin{cases}
  \displaystyle\sup_{x \in \Omega} \lVert  G(x,\cdot)\rVert_{L^p(\Omega;\bar\rho)}^{\gamma_1} & \ell \leq \frac p2 +1 \\
  \displaystyle \sup_{x \in \Omega} \lVert  G(x,\cdot)\rVert_{L^{2p-2(\ell-1)}(\Omega;\bar\rho)}^{\gamma_1} & \ell > \frac p2 +1,
  \end{cases}
\end{align}
which is a bound that is independent of the other variables $x_{i_2},...,x_{i_\ell}$. We proceed in a similar manner for the other terms in the product, i.e. we can apply H\"older's inequality with exponents $\gamma_k/I_p((i,j))$ to obtain
\begin{align}
 &  \int_{\Omega}\left(\prod_{(i,j)\in C_{k,I_p}}|G(x_{i,},x_{j})|^{I_p((i,j))}\right) \, \dx{\bar\rho}(x_{i_k}) \\
\leq & \,  \, \sup_{x \in \Omega} \lVert  G(x,\cdot)\rVert_{L^{\gamma_k}(\Omega;\bar\rho)}^{\alpha_k}\\
 \leq &  \,
 \begin{cases}
  \displaystyle\sup_{x \in \Omega} \lVert  G_{\eps}(x,\cdot)\rVert_{L^p(\Omega;\bar\rho)}^{\gamma_k} & \ell \leq \frac p2 +1 \\
  \displaystyle \sup_{x \in \Omega} \lVert  G_{\eps}(x,\cdot)\rVert_{L^{2p-2(\ell-1)}(\Omega;\bar\rho)}^{\gamma_k} & \ell > \frac p2 +1
 \end{cases}
 \, .
\end{align}
Putting these bounds together we obtain that for any multiindex $I_p\in \mathcal{E}_{p,l}$, we have
\begin{align}
    &\int_{\Omega^N}\prod_{(i,j)\in A}|G(x_{i,},x_{j})|^{I_p((i,j))}\,\mathrm{{d}}\bar\rho^{\otimes N}\\\le & \,  \begin{cases}\displaystyle
  \prod_{k=1}^{\ell-1}\sup_{x \in \Omega} \lVert  G(x,\cdot)\rVert_{L^p(\Omega;\bar\rho)}^{\gamma_k} & \ell \leq \frac p2 +1 \\
  \displaystyle\prod_{k=1}^{\ell-1}\sup_{x \in \Omega} \lVert  G(x,\cdot)\rVert_{L^{2p-2(\ell-1)}(\Omega;\bar\rho)}^{\gamma_k} & \ell > \frac p2 +1
  \end{cases}\\
  =& \,\begin{cases}\displaystyle
  \sup_{x \in \Omega} \lVert  G(x,\cdot)\rVert_{L^p(\Omega;\bar\rho)}^{p} & \ell \leq \frac p2 +1 \\
  \displaystyle\sup_{x \in \Omega} \lVert  G(x,\cdot)\rVert_{L^{2p-2(\ell-1)}(\Omega;\bar\rho)}^{p} & \ell > \frac p2 +1
    \end{cases}\, ,
\end{align}
where we have used the fact that $\sum_{k=1}^{\ell-1}\gamma_k=p$.
Using the above bound, we see that the right hand side of~\eqref{eq:termtosplit} can be estimated as follows
\begin{align}
& \,\frac{{C_p}}{N^{p}}\int_{\Omega^N}\sum_{\ell=2}^p\sum_{I_{p}\in\mathcal{{E}}_{p,\ell}}\prod_{(i,j)\in A}|G(x_{i,},x_{j})|^{I_p((i,j))}\,\mathrm{{d}}\bar\rho^{\otimes N} \\
\leq & \, \frac{C_p}{N^{p}} \sum_{\ell=1}^{\lfloor\frac p2\rfloor +1}|\mathcal{E}_{p,\ell}| \sup_{x \in \Omega} \lVert  G(x,\cdot)\rVert_{L^p(\Omega;\bar\rho)}^{p}\\
& \, + \frac{C_p}{N^{p}} \sum_{\ell=\lceil\frac p2\rceil +1}^{p}|\mathcal{E}_{p,\ell}| \sup_{x \in \Omega} \lVert  G(x,\cdot)\rVert_{L^{2p-2(\ell-1)}(\Omega;\bar\rho)}^{p}\\
\leq & \, \frac{{C_p}}{N^{p}} \sum_{\ell=2}^{\lfloor\frac p2\rfloor +1} N^{\ell}\sup_{x \in \Omega} \lVert  G_{\eps}(x,\cdot)\rVert_{L^p(\Omega;\bar\rho)}^{p}\\
& \, + \frac{{C_p}}{N^{p}} \sum_{\ell=\lceil\frac p2\rceil +1}^{p} N^{\ell} \sup_{x \in \Omega} \lVert  G(x,\cdot)\rVert_{L^{2p-2(\ell-1)}(\Omega;\bar\rho)}^{p}\label{eq:lastineq} \, ,
\end{align}
where we have used~\eqref{eq:sizeE} along with the standard upper bound on the binomial coefficient and we continue to denote the new constant by $C_p$. We now pick some $\gamma \in [\frac12,1)$ and bound the right hand side of~\eqref{eq:lastineq} as follows
\begin{align}
&\,\frac{{C_p}}{N^{p}} \sum_{\ell=2}^{\lfloor\gamma p\rfloor +1}N^{\ell}\sup_{x \in \Omega} \lVert  G(x,\cdot)\rVert_{L^p(\Omega;\bar\rho)}^{p}\, + \frac{{C_p}}{N^{p}} \sum_{\ell=\lceil\gamma p\rceil +1}^{p} N^{\ell} \sup_{x \in \Omega} \lVert  G(x,\cdot)\rVert_{L^{2p-2(\ell-1)}(\Omega;\bar\rho)}^{p} \\
\leq & \,C_p\left( \frac{1}{N^{p-1-\lfloor \gamma p \rfloor}} \sup_{x \in \Omega} \lVert  G(x,\cdot)\rVert_{L^p(\Omega;\bar\rho)}^{p}\, +  \sup_{x \in \Omega} \lVert  G(x,\cdot)\rVert_{L^{2(p-\lceil \gamma p \rceil)}(\Omega;\bar\rho)}^{p}\right) \, ,
\end{align}
where we continue to denote the constant by $C_p$. This completes the proof of the lemma.
\end{proof}

\end{document}
